\providecommand{\sacramentLexiconPageHook}{}
\providecommand{\term}[2]{\textbf{#1} & #2\\\addlinespace[0.12em]}

\begin{landscape}
\sacramentLexiconPageHook
\thispagestyle{plain}
\section*{Metaphysical and Sacramental Lexicon}
\noindent\emph{These terms are analogical instruments of disciplined contemplation. They do not place divine action inside a creaturely mechanism; they distinguish what God causes, what the ecclesial sign signifies, and how a rational creature receives the gift.}

\vspace{0.25em}
\begingroup
\fontsize{7.2}{7.7}\selectfont
\setlength{\tabcolsep}{3pt}
\renewcommand{\arraystretch}{1.01}
\begin{tabularx}{\linewidth}{>{\RaggedRight\arraybackslash}p{0.205\linewidth}>{\RaggedRight\arraybackslash}X}
\toprule
\textbf{Term} & \textbf{Meaning and sacramental use}\\
\midrule
\term{Act / potency}{Potency is capacity to receive perfection; act is possessed perfection. \emph{Obediential potency} names creation's capacity to be elevated by God beyond every naturally due perfection, without making grace a natural claim (Aristotle, \emph{Metaphysics} IX; ST I, q.~105, a.~6).}
\term{Substance / accident}{Substance exists in itself; an accident exists in another. In transubstantiation the whole substance of bread is converted into Christ's Body and the whole substance of wine into his Blood; the sensible accidents remain (Trent XIII.4; ST III, q.~75).}
\term{Matter / sacramental form}{Matter is the determinable sensible element or act; sacramental form is the word or prayer determining its sacred meaning. It is not a substance's intrinsic substantial form; the analogy varies by sacrament (ST III, q.~60, aa.~6--7).}
\term{Remote / proximate}{Remote matter names the material used (water, chrism, oil, bread and wine); proximate matter names its sacramental use (washing, anointing, offering). ``Quasi-matter'' names morally sensible human acts, especially in Penance.}
\term{Sign / thing (\latin{signum / res})}{A sign leads knowledge beyond itself; the thing is what the sign signifies. A sacrament is a sign that Christ also uses to cause the sacred reality signified (Augustine, \emph{De doctrina christiana} II.1; ST III, q.~60, a.~1).}
\term{\latin{Sacramentum tantum}}{The outward rite precisely as sign: washing and words, anointing and words, sacramental species, absolving judgment, imposition of hands, or externally manifested acts of matrimonial consent.}
\term{\latin{Res et sacramentum}}{An intermediate reality that is both effected by the outward sign and signifies a further grace: classically a character, Christ's Eucharistic presence, or the marriage bond. The scheme is not assigned uniformly in every school.}
\term{\latin{Res tantum}}{The grace ultimately signified and caused: incorporation, strengthening, union, reconciliation, healing, ministerial or conjugal grace---all ordered to charity and glory.}
\term{Four causes}{Material cause answers ``from what?''; formal, ``as what?''; efficient, ``by whom?''; final, ``for what end?'' Sacramental analysis uses all four while naming God the transcendent first cause (Aristotle, \emph{Physics} II.3; ST III, qq.~60--65).}
\term{Principal / instrumental}{The Triune God is principal cause; Christ's humanity is the conjoined instrument, and ecclesial rites and ministers are separated instruments. An instrument truly causes by received power (ST III, q.~62, a.~5; q.~64, a.~3).}
\term{Disposition / \latin{obex}}{Disposition prepares the recipient for perfection; an \emph{obex} is an obstacle placed by culpable nonreceptivity. It can impede fruit without defeating the valid sacrament or Christ's promise (ST III, q.~69, aa.~8--10).}
\term{Habit / virtue}{A \latin{habitus} is a stable quality disposing a power to act well or badly. Sanctifying grace and the infused virtues heal and elevate the person for supernatural acts and beatitude (ST I--II, qq.~49--55, 110).}
\term{Sanctifying, actual, and sacramental grace}{Sanctifying grace is stable created participation in divine life; actual grace moves and helps acts; sacramental grace orders such life and helps to each sacrament's proper end (ST III, q.~62, a.~2).}
\term{Character}{An indelible spiritual sign configuring and deputing the recipient to Christian worship. Baptism, Confirmation, and Orders imprint distinct characters and therefore are not repeated (ST III, q.~63; CIC 845).}
\term{\latin{Ex opere operato / operantis}}{The valid sign is efficacious from Christ's work, not the minister's holiness. \emph{Ex opere operantis} names the recipient's graced disposition and cooperation, which govern fruit rather than sacramental power (Trent VII.8; CCC 1128).}
\term{Intention}{The minister must intend at least to perform what the Church performs. Intention makes the bodily deed a human and ecclesial act; it need not include a complete theology or personal holiness (ST III, q.~64, aa.~8--10).}
\term{Validity, liceity, and fruit}{Validity asks whether the sacrament truly occurs; liceity whether lawfully celebrated; fruitfulness whether its grace is unimpeded in this recipient. The three questions must not be collapsed.}
\term{Subject}{The recipient as a determinate, capable person: Baptism precedes every other sacrament; further requirements arise from the nature of each sign and the Church's stewardship (CIC 842--843).}
\term{Participation / adoptive filiation}{A creature possesses by reception and likeness what God is by essence. Grace makes the person God's adopted child and a real participant in divine life without ceasing to be a creature (2~Pet. 1:4; Gal. 4:4--7; ST I, q.~4, a.~3).}
\term{Teleology}{Every nature and action is intelligible through an end (\emph{telos}). The proximate ends differ, but every sacrament is ordered through charity and the Eucharist to resurrection, communion, and the vision of God (ST III, q.~65, a.~3).}
\term{Primary / secondary / ultimate end}{A primary proper end gives an act or institution its governing objective order; secondary ends are real subordinate goods, not disposable extras. The ultimate objective end is God's glory; the recipient's is eternal life and the vision of God.}
\term{Nature / grace}{Grace presupposes, heals, elevates, and perfects nature; it neither treats creation as evil nor makes supernatural beatitude naturally due (ST I, q.~1, a.~8 ad 2; I--II, q.~109).}
\term{Mystery / sacrament}{Greek \emph{mysterion} stresses the divine saving counsel revealed and enacted; Latin \emph{sacramentum} stresses consecrating sign and bond. East and West accent the same Christological economy from complementary vocabularies.}
\bottomrule
\end{tabularx}
\endgroup
\end{landscape}
